% Generated draft: requires mathematical and source review.
\documentclass[11pt]{article}
\usepackage[margin=25mm]{geometry}
\usepackage{fontspec}
\setmainfont{Latin Modern Roman}
\usepackage{amsmath,amssymb,mathtools,unicode-math}
\setmathfont{Latin Modern Math}
\usepackage{enumitem,needspace,xurl,hyperref}
\setlength{\parindent}{0pt}
\setlength{\parskip}{5pt}
\setlength{\emergencystretch}{3em}
\newcommand{\N}{\mathbb{N}}
\newcommand{\Z}{\mathbb{Z}}
\newcommand{\Q}{\mathbb{Q}}
\newcommand{\R}{\mathbb{R}}
\newcommand{\C}{\mathbb{C}}
\newcommand{\F}{\mathbb{F}}
\newcommand{\A}{\mathbb{A}}
\newcommand{\PP}{\mathbb{P}}
\newcommand{\E}{\mathbb{E}}
\newcommand{\eps}{\varepsilon}
\newcommand{\dd}{\,\mathrm{d}}
\begin{document}
\section*{UnsolvedMath ID: 701599}
\textbf{Source problem code:} catalogue entry
\subsection{Definitions and mathematical statement}
Throughout, groups are finite. We use the conventional interpretation of \emph{irreducible character} as an ordinary irreducible complex character; the supplied source does not explicitly specify the character field. Write \(\operatorname{Irr}(G)\) for these characters. A linear character of a subgroup \(H\leq G\) is a character of a one-dimensional complex representation of \(H\). The group \(G\) is \emph{monomial} if, for every \(\chi\in\operatorname{Irr}(G)\), there are a subgroup \(H\leq G\) and a linear character \(\lambda\) of \(H\) such that
\[
\chi=\operatorname{Ind}_H^G\lambda.
\]
It is \emph{normally monomial} if \(H\) can always be chosen normal in \(G\), and \emph{subnormally monomial} if \(H\) can always be chosen subnormal in \(G\). Here subnormal means that there is a finite chain \(H=H_0\trianglelefteq H_1\trianglelefteq\cdots\trianglelefteq H_r=G\). The subgroup and linear character may depend on \(\chi\).

Define \(G^{(0)}=G\) and \(G^{(i+1)}=[G^{(i)},G^{(i)}]\). The group is soluble if \(G^{(d)}=1\) for some \(d\geq0\); its \emph{derived length} \(d(G)\) is the least such \(d\). The supplied source attributes to Taketa the assertion that finite monomial groups are soluble, so derived length is finite for the groups considered here.

The source does not define \emph{nilpotence length}. For part (b), we interpret it in the conventional finite soluble group sense, namely the \emph{Fitting height} \(h(G)\): the least integer \(r\geq0\) for which there is a series
\[
1=N_0\leq N_1\leq\cdots\leq N_r=G,
\qquad N_i\trianglelefteq G,
\]
with every factor \(N_i/N_{i-1}\) nilpotent. A group \(K\) is nilpotent if its lower central series, defined by \(\gamma_1(K)=K\) and \(\gamma_{i+1}(K)=[\gamma_i(K),K]\), reaches \(1\). The intended meaning of nilpotence length requires source confirmation; part (b) below is conditional on this interpretation.

\begin{enumerate}
\item[(a)] Is it true that for every prime \(p\) and every integer \(m\geq1\), there exists a normally monomial finite \(p\)-group \(G\) such that
\[
d(G)\geq m?
\]
Here a finite \(p\)-group has order \(p^n\) for some integer \(n\geq0\), and \(p\) is fixed while \(m\) varies.
\item[(b)] Is it true that for every integer \(m\geq1\), there exists a subnormally monomial finite group \(G\) such that
\[
h(G)\geq m?
\]
\end{enumerate}

For context, the supplied source states that metabelian groups, meaning groups with \(G^{(2)}=1\), are normally monomial, and that abelian-by-nilpotent groups are subnormally monomial. Here abelian-by-nilpotent means that \(G\) has an abelian normal subgroup \(A\) with nilpotent quotient \(G/A\). It also states that normally monomial finite groups of arbitrarily large derived length exist, without imposing the fixed-prime \(p\)-group restriction of part (a). These background assertions are retained as supplied source claims.
\subsection{Short English statement}
For each fixed prime \(p\), can normally monomial finite \(p\)-groups have unbounded derived length? Can subnormally monomial finite groups have unbounded nilpotence length, interpreted here as Fitting height? Normal or subnormal monomiality requires every irreducible complex character to be induced from a linear character of a normal or subnormal subgroup, respectively. The supplied registry labels the problem open; this draft does not independently verify that status.
\subsection{Sources}
\begin{itemize}
\item Local UnsolvedMath registry: supplied complete record with permanent ID \(701599\), exact source problem code \texttt{catalogue entry}, and title \emph{Kourovka Notebook Problem 19.51}. The supplied status is \texttt{open}; it is a registry label, not an independently verified mathematical assessment.
\item Supplied source attribution: \emph{The Kourovka Notebook (Unsolved Problems in Group Theory)}, 19th issue (2018), Problem 19.51; proposed by A. Mann in 2018. The background also identifies the 21st issue edition (2026), \texttt{arXiv:1401.0300v48}. These bibliographic details are supplied and have not been independently checked.
\item The supplied collection description names E. I. Khukhro and V. D. Mazurov as editors of the Kourovka Notebook collection covering Issues 1--21 (1965--2026).
\item The source attributes solubility of finite monomial groups to Taketa, but supplies no separate bibliographic citation for that result or for the other background assertions.
\end{itemize}
\end{document}
